%----------------------------------------------------------------------------------------
% VEELTERMFUNCTIES. Alle opgaven zijn eigen werk.
%
% Opbouw: Start, leerstof (theorie, werkwijze of uitleg, voorbeeld, korte oefening), en op een
% nieuwe pagina de Opdrachten (door elkaar, stijgende moeilijkheid), Koppel en Test jezelf.
% Labels van dit hoofdstuk beginnen met h1-.
%----------------------------------------------------------------------------------------

\mychapter{1}{Veeltermfuncties}

\lkr{Leerplan B+S': LPD 18 en K1 (rationale functies), met de veeltermfuncties als voorbereiding: delen, de reststelling, de factorstelling en het schema van Horner, ontbinden en de grafiek schetsen vanuit nulwaarden, multipliciteit en de hoogstegraadsterm. Eerste- en tweedegraadsfuncties kennen leerlingen uit de 2de graad.}

Eerste- en tweedegraadsfuncties ken je al: rechten en parabolen. Een veeltermfunctie van hogere graad kan meer bochten maken. Om haar grafiek te schetsen heb je vooral haar nulwaarden nodig, en daarvoor moet je de veelterm ontbinden. Dat doe je door te delen.

\begin{oefeningen}[Start]
\oef{h1-start-1}

\oef{h1-verken}
\end{oefeningen}

%========================================================================================
\section{Veeltermen delen}

\begin{theorie}
Een \textbf{\textcolor{accent}{veeltermfunctie}} van graad $n$ is $f(x) = a_n x^n + a_{n-1}x^{n-1} + \cdots + a_1 x + a_0$ met $a_n \neq 0$.

\begin{stelling}{Euclidische deling}
Voor veeltermen $A(x)$ en $B(x) \neq 0$ bestaan er unieke veeltermen $Q(x)$ (het quotiënt) en $R(x)$ (de rest) met
\[\formule{A(x) = B(x) \cdot Q(x) + R(x)} \qquad \text{en} \qquad \formule{\text{graad } R < \text{graad } B}.\]
\end{stelling}
\begin{stelling}{Reststelling}
De rest bij deling van $A(x)$ door $x - a$ is $\formule{A(a)}$.
\end{stelling}
\begin{bewijs}
De rest heeft graad kleiner dan 1: een getal $r$. Uit $A(x) = (x - a)Q(x) + r$ volgt voor $x = a$:
\bewijsregel{A(a) = 0 \cdot Q(a) + r = r.}
\end{bewijs}
\end{theorie}

Delen door $x - a$ gaat snel met het \textbf{schema van Horner}: zet de coëfficiënten van $A(x)$ op een rij (vergeet geen nullen), breng de eerste naar beneden, en herhaal: vermenigvuldig met $a$, zet het resultaat in de volgende kolom en tel op.

\voorbeeld{}{Deel $A(x) = 2x^3 - 3x^2 + 4x - 5$ door $x - 2$.\label{vb:h1-horner}}{%
\naast{%
\uitleg{Hier is $a = 2$. Het laatste getal van de onderste rij is de rest, de andere zijn de coëfficiënten van het quotiënt (een graad lager).}
\begin{stappenvrij}{4.6cm}{3.6cm}
Q(x) = 2x^2 + x + 6 & Coëfficiënten $2, 1, 6$. \\
R = 7 & Het laatste getal. \\
A(2) = 16 - 12 + 8 - 5 = 7 & Controle met de reststelling.
\end{stappenvrij}}{%
{\renewcommand{\arraystretch}{1.3}%
$\begin{array}{c|rrrr}
 & 2 & -3 & 4 & -5 \\
2 & & 4 & 2 & 12 \\ \hline
 & 2 & 1 & 6 & \textcolor{accent}{\mathbf{7}}
\end{array}$}}}

\begin{oefeningen}
\needspace{6\baselineskip}
\oef{h1-horner}
\end{oefeningen}

Door een veelterm van hogere graad deel je met een staartdeling, net zoals bij getallen.

\voorbeeld{}{Deel $A(x) = x^4 + 2x^3 - x + 3$ door $B(x) = x^2 + 1$.\label{vb:h1-staartdeling}}{%
\uitleg{Deel telkens de hoogstegraadsterm van wat overblijft door $x^2$, vermenigvuldig terug met $B(x)$ en trek af. Stop als de graad kleiner is dan 2.}
\begin{stappenbreed}
x^4 + 2x^3 - x + 3 - x^2(x^2 + 1) = 2x^3 - x^2 - x + 3 & Eerste term van $Q$: $x^2$. \\
2x^3 - x^2 - x + 3 - 2x(x^2 + 1) = -x^2 - 3x + 3 & Tweede term: $2x$. \\
-x^2 - 3x + 3 - (-1)(x^2 + 1) = -3x + 4 & Derde term: $-1$. Graad 1: stop. \\
Q(x) = x^2 + 2x - 1,\quad R(x) = -3x + 4 &
\end{stappenbreed}}

\begin{oefeningen}
\oef{h1-delen}

\needspace{6\baselineskip}
\oef{h1-rest}
\end{oefeningen}

%========================================================================================
\section{Ontbinden en nulwaarden}

\begin{theorie}
\begin{stelling}{Factorstelling}
$x - a$ is een deler van $A(x)$ \quad $\formule{\iff}$ \quad $\formule{A(a) = 0}$.
\end{stelling}
Is $A(x) = (x - a)^m \cdot Q(x)$ met $Q(a) \neq 0$, dan is $a$ een nulwaarde met \textbf{\textcolor{accent}{multipliciteit}} $m$.
\end{theorie}

De factorstelling volgt uit de reststelling: de rest is nul als en slechts als $A(a) = 0$. Heeft een veelterm gehele coëfficiënten, dan is elke gehele nulwaarde een deler van de constante term $a_0$. Zo vind je kandidaten om te proberen.

\voorbeeld{}{Ontbind $f(x) = x^3 - 2x^2 - 5x + 6$ en geef de nulwaarden.\label{vb:h1-ontbinden}}{%
\naast{%
\uitleg{Probeer delers van $6$. Vind je een nulwaarde, deel dan met Horner en ontbind het quotiënt (een tweedegraadsveelterm) verder.}
\begin{stappenvrij}{4.9cm}{3.2cm}
f(1) = 1 - 2 - 5 + 6 = 0 & $x - 1$ is een deler. \\
f(x) = (x - 1)(x^2 - x - 6) & Horner. \\
= (x - 1)(x - 3)(x + 2) & Som $1$, product $-6$. \\
\text{nulwaarden: } -2,\ 1,\ 3 &
\end{stappenvrij}}{%
{\renewcommand{\arraystretch}{1.3}%
$\begin{array}{c|rrrr}
 & 1 & -2 & -5 & 6 \\
1 & & 1 & -1 & -6 \\ \hline
 & 1 & -1 & -6 & \textcolor{accent}{\mathbf{0}}
\end{array}$}}}

\begin{werkwijze}
Zonder eerst een gemeenschappelijke factor af en gebruik merkwaardige producten. Probeer daarna de delers van de constante term (positief en negatief) tot je een nulwaarde $a$ vindt. Deel door $x - a$ met Horner, en ontbind het quotiënt verder. Een tweedegraadsveelterm ontbind je met de discriminant.
\end{werkwijze}

\begin{oefeningen}
\oef{h1-ontbinden}

\needspace{6\baselineskip}
\oef{h1-factor}

\needspace{6\baselineskip}
\oef{h1-multipliciteit}
\end{oefeningen}

%========================================================================================
\section{De grafiek van een veeltermfunctie}

Voor grote $|x|$ overheerst de hoogstegraadsterm: in $f(x) = x^3 - 2x^2 - 5x + 6$ is voor $x = 100$ de term $x^3 = 1\,000\,000$, en de rest samen maar ongeveer $-20\,500$. De grafiek gedraagt zich dan zoals die van $a_n x^n$.

\begin{theorie}
De grafiek van een veeltermfunctie
\begin{itemize}
\item gedraagt zich voor grote $|x|$ zoals die van de \textbf{\textcolor{accent}{hoogstegraadsterm}} $a_n x^n$;
\item \textbf{\textcolor{accent}{snijdt}} de $x$-as in een nulwaarde met oneven multipliciteit (het teken wisselt);
\item \textbf{\textcolor{accent}{raakt}} de $x$-as in een nulwaarde met even multipliciteit (het teken blijft).
\end{itemize}
\end{theorie}

\voorbeeld{}{Schets de grafiek van $f(x) = -(x + 1)^2(x - 2)$.\label{vb:h1-schets}}{%
\naast{%
\uitleg{Lees de nulwaarden en hun multipliciteit af, en het gedrag aan de uiteinden. Bereken het snijpunt met de $y$-as.}
\begin{stappenvrij}{4.5cm}{3.3cm}
\text{hoogstegraadsterm } -x^3 & Links boven, rechts onder. \\
-1 \text{ (dubbel)},\ 2 \text{ (enkel)} & Raakt in $-1$, snijdt in $2$. \\
f(0) = -1 \cdot (-2) = 2 & Snijpunt met de $y$-as.
\end{stappenvrij}
\figuren{\tekentabel{cccccc}{x & & -1 & & 2 & }{f(x) & + & 0 & + & 0 & -}}}{%
\begin{tikzpicture}[x=.6cm,y=.42cm,font=\scriptsize]
  \rooster{-3}{1}{3}{-4}{1}{5}
  \kromme{-2.1}{2.3}{-(\x+1)^2*(\x-2)}
  \vol{(-1,0)} \vol{(2,0)} \vol{(0,2)}
\end{tikzpicture}}}

\begin{werkwijze}
Ontbind en bepaal de nulwaarden met hun multipliciteit. Stel de tekentabel op. Bepaal het gedrag aan de uiteinden met de hoogstegraadsterm en bereken $f(0)$. Teken: snijden bij oneven, raken bij even multipliciteit. Waar precies de toppen liggen, bereken je later met afgeleiden.
\end{werkwijze}

\begin{oefeningen}
\oef{h1-schets}

\needspace{6\baselineskip}
\oef{h1-uiteinden}
\end{oefeningen}

%========================================================================================
\opdrachtenkop

\begin{oefeningen}
\needspace{12\baselineskip}
\oef{h1-o1}

\oef{h1-o13}

\needspace{12\baselineskip}
\oef{h1-o3}

\oef{h1-o10}

\needspace{12\baselineskip}
\oef{h1-o15}

\oef{h1-o2}

\needspace{8\baselineskip}
\oef{h1-o12}

\needspace{8\baselineskip}
\oef{h1-o14}

\oef{h1-o11}

\oef{h1-o4}

\needspace{8\baselineskip}
\oef{h1-o17}

\needspace{14\baselineskip}
\oef{h1-o19}

\needspace{9\baselineskip}
\oef{h1-o20}

\oef{h1-o6}

\needspace{8\baselineskip}
\oef{h1-o16}

\needspace{10\baselineskip}
\oef{h1-o23}

\oef{h1-o7}

\needspace{8\baselineskip}
\oef{h1-o21}

\oef{h1-o5}

\needspace{8\baselineskip}
\oef{h1-o22}

\needspace{14\baselineskip}
\oef{h1-o18}

\oef{h1-o8}

\needspace{8\baselineskip}
\oef{h1-o24}

\needspace{10\baselineskip}
\oef{h1-o25}

\needspace{8\baselineskip}
\oef{h1-o27}

\needspace{8\baselineskip}
\oef{h1-o26}

\needspace{6\baselineskip}
\oef{h1-o9}
\end{oefeningen}

%----------------------------------------------------------------------------------------
% KOPPEL: verbanden met andere leerstof
%----------------------------------------------------------------------------------------
\needspace{20\baselineskip}
\opdrachtsectie{Koppel}

\lkr{Complexe getallen en Algemene functies: kijk na wat al gezien is. Het pakket legt de volgorde niet vast.}

\begin{oefeningen}
\oef{h1-koppel-1}

\oef{h1-koppel-3}
\end{oefeningen}

%----------------------------------------------------------------------------------------
\opdrachtsectie{Test jezelf}

\begin{oefeningen}
\oef{h1-t1}

\oef{h1-t2}

\oef{h1-t5}

\oef{h1-t3}

\oef{h1-t4}

\oef{h1-t6}
\end{oefeningen}

\herhaalschema{%
  \fv{oef:h1-t1}{vb:h1-horner}{oef:h1-horner}%
  \fv{oef:h1-t2}{vb:h1-staartdeling}{oef:h1-delen}%
  \fv{oef:h1-t5}{vb:h1-ontbinden}{oef:h1-factor}%
  \fv{oef:h1-t3}{vb:h1-ontbinden}{oef:h1-ontbinden}%
  \fv{oef:h1-t4}{vb:h1-schets}{oef:h1-schets}%
  \fv{oef:h1-t6}{vb:h1-schets}{oef:h1-o6}}

\ifzelfstudie
\section*{Maak je eigen samenvatting}
Sluit dit boek en schrijf in je schrift: de euclidische deling, het schema van Horner, de reststelling met bewijs, de factorstelling, hoe je kandidaat-nulwaarden vindt, wat multipliciteit betekent voor de grafiek, en de werkwijze om een veeltermfunctie te schetsen. Vergelijk daarna met de kaders.
\fi
